\subsection{Paměťové nároky}
\label{subsec:llm-memory-requirements}

Paměťové nároky pro provoz jazykového modelu jsou výrazně vyšší, než by se mohlo na první pohled zdát.
Celkové nároky tvoří tři složky:

\begin{itemize}
    \item paměť potřebnou pro uložení vah modelu,
    \item paměť využívanou během inference (mezivýpočty a attention cache),
    \item režii spojenou s paralelním zpracováním více požadavků.
\end{itemize}

Tyto faktory omezují velikost modelu, který lze efektivně provozovat lokálně, a zároveň určují maximální počet paralelních analýz bez výrazného zpomalení systému.

\subsubsection*{Váhy modelu}

V plné 16-bitové přesnosti (FP16 nebo BF16) jsou potřeba 2 bajty na parametr.
Model se 7 miliardami parametrů tak vyžaduje přibližně 14 GB VRAM pouze pro uložení vah.
Model se 13 miliardami parametrů potřebuje přibližně 26 GB\@.
Tyto hodnoty přesahují kapacitu většiny běžných spotřebitelských grafických karet.

\textbf{Kvantizace} je technikou, která tuto situaci výrazně mění.
Redukcí přesnosti vah z 16 bitů na 8 nebo 4 bity lze paměťové nároky snížit na přibližně polovinu, respektive čtvrtinu původní hodnoty.
Model 7B v 4-bitové kvantizaci (označované například jako GGUF Q4\_K\_M) vyžaduje přibližně 4 až 5 GB VRAM, hodnota dostupná i na středně výkonných spotřebitelských GPU\@.
Ztráta kvality výstupu při kvalitní kvantizaci je pro praktické použití zpravidla zanedbatelná a pro účely analýzy studentských řešení zcela přijatelná.

\begin{table}[H]
    \centering
    \resizebox{\textwidth}{!}{%
        \begin{tabular}{lllll}
            \toprule
            \textbf{Velikost modelu} & \textbf{FP16} & \textbf{Q8} & \textbf{Q4} & \textbf{Typická GPU} \\
            \midrule
            3B parametrů   & $\sim$6 GB   & $\sim$3 GB   & $\sim$2 GB   & RTX 4060 (8 GB) \\
            7B parametrů   & $\sim$14 GB  & $\sim$7 GB   & $\sim$4 GB   & RTX 3060 (12 GB) \\
            13B parametrů  & $\sim$26 GB  & $\sim$13 GB  & $\sim$7 GB   & RTX 4070 Ti (12 GB) \\
            34B parametrů  & $\sim$68 GB  & $\sim$34 GB  & $\sim$19 GB  & RTX 4090 (24 GB) \\
            70B parametrů  & $\sim$140 GB & $\sim$70 GB  & $\sim$38 GB  & 2$\times$ RTX 4090 (48 GB) \\
            \bottomrule
        \end{tabular}
    }
    \caption{Přibližné paměťové nároky modelů podle počtu parametrů a kvantizace (bez KV cache)}
    \label{tab:memory-requirements}
\end{table}

\subsubsection*{KV cache}
\label{subsubsec:kv-cache}

KV cache byla vysvětlena v \secref[sekci]{subsec:inference-generation}.
Její velikost roste lineárně s délkou zpracovávané sekvence a s počtem vrstev modelu.
Pro model se 32 vrstvami, kontextovým oknem 8~000 tokenů a jedním souběžným požadavkem je KV cache typicky v řádu 2 až 4 GB VRAM, v závislosti na architektuře attention mechanismu
(MHA -- Multi-Head Attention vs.\ GQA -- Grouped-Query Attention).
Při zpracování více souběžných požadavků (dávkové zpracování) se tento nárok násobí.

\subsubsection*{Praktické dopady}
\label{subsubsec:practical-memory-implications}

Z výše uvedených hodnot plyne, že pro efektivní lokální provoz modelu třídy 7B parametrů je zapotřebí GPU s alespoň 4 GB VRAM pro Q4 kvantizaci, přičemž 12 GB nebo více poskytuje pohodlnou rezervu pro KV cache.
Model třídy 13B parametrů vyžaduje GPU s 8 až 12 GB VRAM v Q4 kvantizaci.
Pro modely třídy 70B a více parametrů je nutné nasazení více GPU -- například konfigurace 2$\times$ RTX 4090 s celkovou kapacitou 48 GB VRAM umožňuje provoz 70B modelu v Q4 kvantizaci.

%Pokud škola disponuje odpovídajícím GPU (například NVIDIA RTX 3060 s 24 GB nebo profesionální kartou řady A), je lokální nasazení modelu 7B nebo 13B technicky realistické.
%V opačném případě je variantou cloudové API, u nějž jsou paměťové nároky plně v režiji poskytovatele.

\endinput